%************************************************
\chapter*{Coda: The Next Rung}\phantomsection\label{ch:coda}
\addcontentsline{toc}{chapter}{\tocEntry{Coda: The Next Rung}}
\markboth{\spacedlowsmallcaps{Coda: The Next Rung}}{\spacedlowsmallcaps{Coda: The Next Rung}}
%************************************************

\chapterlead{Where the new numbers stand on the long ladder of numbers, what they still need, and the question this book hands to its readers.}

Pick up the pebble again, and let it go. Rub the balloon on your sleeve, and hold it near your hair. At the start of this book these were two pulls, described by two laws, and every particle carried two numbers that the two laws read separately. Now they can be written with one number for each particle and one law for each pair, and, once the bodies move, with one set of fields for both. In the new language, the pebble's fall belongs to the first term of the Cosmic Interaction Law, mass with mass, and your rising hair to the second, charge with charge. Between them the law holds a third term, mass with charge. Read by its real part, as Chapter~\ref{ch:interaction} reads every force we measure, that term adds nothing to either pull.

%------------------------------------------------
\section*{The ladder, one rung higher}\phantomsection
\addcontentsline{toc}{section}{The ladder, one rung higher}
\markright{\spacedlowsmallcaps{The ladder, one rung higher}}
%------------------------------------------------

Chapter~\ref{ch:one-number} set out the ladder of numbers as far as the complex numbers (\Cref{tab:ladder-small}). The ladder did not stop there. The table below continues it with the numbers that came after $\C$, and with the rung proposed in this book. Each row records a question that the numbers of the time could not answer, and what the new numbers later allowed physics to say \citep{Kline1972,Baez2002}. A map of the whole road, from the first notched bones to the present day, is given in \Cref{app:long-road}.

\begin{center}
\small
\begin{tabularx}{\textwidth}{@{}>{\raggedright\arraybackslash}p{0.25\textwidth} >{\raggedright\arraybackslash}X >{\raggedright\arraybackslash}X@{}}
\toprule
\tableheadline{New numbers} & \tableheadline{The question} & \tableheadline{What followed} \\
\midrule
negative numbers & What is $3 - 5$? & direction, deficit, opposite charges \\
irrational numbers & How long is the diagonal of a unit square? & continuous magnitude, the real line \\
complex numbers $\C$ & Which number squares to $-1$? & rotation, phase, alternating currents, quantum amplitudes \\
quaternions $\mathbb{H}$, 1843 & Can the arithmetic of the plane be carried into space? & the motion of rigid bodies, spin one half \\
Grassmann numbers, 1844 & Can a product change sign when its order changes, $ab = -ba$? & the particles of matter in quantum theory; supersymmetry, not yet observed \\
Clifford algebras, 1878 & Can the algebras of Hamilton and Grassmann be joined into one? & the Dirac equation of the electron, antimatter \\
octonions $\mathbb{O}$, 1843 to 1845 & Can the quaternions be extended once more? & speculative: the symmetries of particle physics \\
\midrule
General Numbers $\GN$, this book & Must the second axis of the number plane stand at a right angle? & proposed: one number for mass and charge \\
\bottomrule
\end{tabularx}
\end{center}

Two of these rungs appeared earlier in this book. Hamilton's quaternions of 1843, cut into the stone of a Dublin bridge, opened Chapter~\ref{ch:general-numbers} \citep{Hamilton1844}. Grassmann's numbers, whose products change sign when their order changes, return below \citep{Grassmann1844}. The Clifford algebras \citep{Clifford1878} had to wait fifty years for Paul Dirac, who needed a square root of a differential operator, was led to matrices that anticommute, and found the equation that gave the electron its spin and predicted antimatter \citep{Dirac1928}.

%------------------------------------------------
\section*{What every rung needed}\phantomsection
\addcontentsline{toc}{section}{What every rung needed}
\markright{\spacedlowsmallcaps{What every rung needed}}
%------------------------------------------------

Look down the ladder and a pattern appears. It is worth stating plainly, because it tells us what the General Numbers still need. Every rung needed four things.
\begin{enumerate}
  \item \emph{A question} that the old numbers could not answer, sharp enough that no amount of ingenuity inside the old system could dissolve it.
  \item \emph{A definition} that contradicted nothing, and that kept the old numbers inside the new ones as a special case.
  \item \emph{Consequences} that nobody had put in on purpose: the surprises that made the new numbers worth having.
  \item In the end, \emph{a fact about the world} that could not be said without them.
\end{enumerate}

The General Numbers, as this book presents them, offer a candidate for each of the first three. The question is the one Chapter~\ref{ch:one-number} asked: can a single number carry both the mass and the charge of a particle? Here honesty requires a remark. A complex number, such as $m + iq\sqrt{K/G}$, can also hold a mass and a charge in one number, and each relative plane of the General Number Plane is, algebraically, the complex plane drawn in a leaning frame. What the General Numbers add is the frame itself: the angle between the two axes, and the length of the second unit. Whether nature needs that freedom, as it turned out to need the number $i$, is part of what remains to be shown. The definition is Axiom~\ref{ax:gn}, which contains the complex numbers as the relative plane $\pb = \pi/2$ (\Cref{thm:classical-recovery}). The consequences include the spiral of powers, the triangle of every general number, the new identity $(r, \alpha, \phi)$ of a particle, and the Rule of the Square, whose minus sign pushes like charges apart and, in Chapter~\ref{ch:unified}, gives the energy of the charge fields a plus sign and the energy of gravity a minus sign.

The fourth step is not in any book. For the complex numbers it took nearly three and a half centuries, from Cardano's puzzle of 1545 to Steinmetz's alternating currents of 1893, and longer still to quantum mechanics, where complex numbers turned out to be not merely convenient but necessary. For the General Numbers, the fourth step will be written, if it is written at all, by an experiment.

%------------------------------------------------
\section*{Earlier dreams of one number}\phantomsection
\addcontentsline{toc}{section}{Earlier dreams of one number}
\markright{\spacedlowsmallcaps{Earlier dreams of one number}}
%------------------------------------------------

The wish to hold mass and charge in a single mathematical object is not new, and it is only fair to place this book beside the attempts that came before it.

In 1921 Theodor Kaluza wrote Einstein's theory of gravity in five dimensions instead of four, and found that the extra parts of the geometry behave exactly like the electromagnetic field \citep{Kaluza1921}. In 1926 Oskar Klein added that the fifth dimension could be curled up into a tiny circle. Motion around that circle then comes in whole steps, and those steps behave like electric charge \citep{Klein1926}. In the theory of Kaluza and Klein, charge became motion in a direction we cannot see, and mass and charge became two parts of one geometric object. The theory ran into physical difficulties and did not become the physics of our world, although its central idea still lives in modern research.

In 1939 Eugene Wigner showed that the mass of a particle is a label: it tells how the particle behaves under the symmetries of space and time \citep{Wigner1939}. Electric charge, it became clear, labels a symmetry of a different kind, an internal one, attached to each point of space without being part of its geometry. In 1967 Sidney Coleman and Jeffrey Mandula proved that, under very general assumptions, these two kinds of symmetry cannot be mixed in any straightforward way \citep{ColemanMandula1967}. But every good impossibility theorem is also a map of its own loophole. In 1975 Rudolf Haag, Jan {\L}opusza{\'n}ski and Martin Sohnius showed exactly where this one lies: the mixing becomes possible, and only then, if some quantities multiply with a change of sign, $ab = -ba$, as Grassmann's numbers of 1844 do \citep{HLS1975}. That loophole is supersymmetry, which physicists had begun to explore a few years earlier. So far, it has not been seen in any experiment.

Two lessons follow. The first is that the question of this book is a serious one: some of the best physicists of the twentieth century asked it too. The second is that every serious attempt to answer it has needed a new mathematical structure: an extra dimension of space, or a new kind of number. The General Numbers are a different and far more elementary proposal. They add no dimension to space and no new symmetry to physics. They change only the frame of the number plane, by letting its second axis lean. Whether so small a change can carry so large a load is exactly what remains to be seen.

%------------------------------------------------
\section*{The open question}\phantomsection
\addcontentsline{toc}{section}{The open question}
\markright{\spacedlowsmallcaps{The open question}}
%------------------------------------------------

Faraday asked with his falling coil whether there is a link between mass and charge beyond the separate laws of Newton and Coulomb. The Cosmic Mass Theory gives an answer, though not the one he looked for. Its law has a third term, the \emph{Coupling Force}, which pairs the mass of each body with the charge of the other. Read by its real part, as Chapter~\ref{ch:interaction} reads every measured force, that term adds nothing along the line between two bodies, and Faraday's needle, which stayed at zero, is not asked to move. The link is the law itself: one number for each body, one product for gravity and electricity, and, once the bodies move, one set of fields for both (Chapter~\ref{ch:unified}).

If this book leaves one question in the reader's hands, it is what that link means in the world, and it is the question that Chapter~\ref{ch:standard-model} leaves open at its end: the part of a cosmic mass that lies across the mass axis, and the Composer coordinate of a product, which holds the Coupling term together with a share of the term of two charges, far larger for charged particles (\Cref{subsec:law-read}), but is not read as a force. The Composer axis itself is held at $79.86^\circ$, the angle at which the plane shows why like charges push apart (\Cref{subsec:which-root}). Volume~2 of this work takes up the Composer coordinate, beginning with the Coupling term (\Cref{sec:al-coupling}). Chapter~\ref{ch:open} gathers the readings proposed for it, with what each would require and what measurement already says. Chapter~\ref{ch:allows} draws the whole frame: antimatter, on the same lines as matter; the half of negative mass, empty so far; and the black holes, far out on the mass axis.

Whether it stands for something in the world is not for the author to decide, nor for the reader. It is for the world. The way forward is the way every rung of the ladder was made secure: a sharp prediction, and a measurement able to confirm it or to refute it. The first chapter of this book promised to state every step clearly, so that you could follow it and judge it for yourself. If you have come this far, you are in a position to do exactly that, and perhaps to design the experiment.

%------------------------------------------------
\section*{The characters of the book}\phantomsection
\addcontentsline{toc}{section}{The characters of the book}
\markright{\spacedlowsmallcaps{The characters of the book}}
%------------------------------------------------

Galileo wrote that the book of nature is written in the language of mathematics, and that its characters are triangles, circles and other geometric figures. He did not say right triangles. For four centuries physics has read that book with an alphabet built on the right angle: the right triangle behind the sine and the cosine, the perpendicular axes of the complex plane. This book has tried to add one letter to that alphabet: a triangle that leans, and a number plane that leans with it.

Again and again, as the ladder shows, somebody has written down an object that seemed to have no right to exist, and found that it worked. There is no reason to think it has happened for the last time.
